%% paper65_graphentheorie_de.tex
%% Graphentheorie: Grundlagen, Färbung, Planarität und das Hadwiger-Nelson-Problem
%% Autor: Michael Fuhrmann
%% Projekt: Specialist Mathematics Project, Build 145
%% Datum: 2026-03-12

\documentclass[reqno,12pt]{amsart}

%% ---------------------------------------------------------------
%% Pakete
%% ---------------------------------------------------------------
\usepackage[ngerman]{babel}
\usepackage{amsmath, amssymb, amsthm}
\usepackage{mathtools}
\usepackage{geometry}
\usepackage{hyperref}
\usepackage{enumitem}
\usepackage{tikz}
\usepackage{xcolor}
\usepackage{booktabs}
\usepackage{microtype}

\geometry{margin=2.5cm}

\hypersetup{
  colorlinks=true,
  linkcolor=blue!60!black,
  citecolor=green!50!black,
  urlcolor=blue!70!black
}

%% ---------------------------------------------------------------
%% Theorem-Umgebungen (Deutsch)
%% ---------------------------------------------------------------
\newtheorem{theorem}{Satz}[section]
\newtheorem{lemma}[theorem]{Lemma}
\newtheorem{corollary}[theorem]{Korollar}
\newtheorem{proposition}[theorem]{Proposition}
\theoremstyle{definition}
\newtheorem{definition}[theorem]{Definition}
\newtheorem{example}[theorem]{Beispiel}
\theoremstyle{remark}
\newtheorem{remark}[theorem]{Bemerkung}
\newtheorem{conjecture}[theorem]{Vermutung}

%% ---------------------------------------------------------------
%% Eigene Befehle
%% ---------------------------------------------------------------
\newcommand{\Z}{\mathbb{Z}}
\newcommand{\R}{\mathbb{R}}
\newcommand{\N}{\mathbb{N}}
\newcommand{\C}{\mathbb{C}}
\newcommand{\chr}{\chi}
\newcommand{\grad}{\mathrm{grad}}

%% ---------------------------------------------------------------
%% Dokument
%% ---------------------------------------------------------------

\title{Graphentheorie: Grundlagen, Färbung, Planarität\\
       und das Hadwiger-Nelson-Problem}

\author{Michael Fuhrmann}

\address{Specialist Mathematics Project, Build 145}

\email{specialist-maths@localhost}

\date{2026-03-12}

\keywords{Graphentheorie, Graphenfärbung, planare Graphen, Spannbäume,
          spektrale Graphentheorie, Hadwiger-Nelson-Problem, Vierfarbensatz,
          chromatische Zahl, Euler-Formel, Menger-Satz}

\subjclass[2020]{05C15, 05C10, 05C40, 05C50, 05C45, 05D10}

\begin{abstract}
Die Graphentheorie, entstanden aus Eulers Lösung des Königsberger Brückenproblems
im Jahr 1736, hat sich zu einem der vielseitigsten Gebiete der diskreten
Mathematik entwickelt, mit Anwendungen in Informatik, Physik, Chemie und
Kombinatorik. Diese Arbeit präsentiert eine in sich geschlossene Entwicklung
der zentralen Ergebnisse der Graphentheorie: das Handschlag-Lemma,
Charakterisierungen von Bäumen (einschließlich der Cayley-Formel
$\kappa(K_n) = n^{n-2}$), der Menger-Satz und der Max-Flow-Min-Cut-Satz,
Graphenfärbetheorie mit dem berühmten Vierfarbensatz (Appel-Haken 1976;
Robertson et al.\ 1997), Planarität und Kuratowskis Charakterisierung durch
$K_5$ und $K_{3,3}$, die Euler-Formel $V - E + F = 2$, der Fünffarbensatz,
spektrale Graphentheorie mit der Cheeger-Ungleichung, die Theorie Eulerscher
und Hamiltonischer Graphen sowie die NP-Vollständigkeit des
Hamiltonischen-Pfad-Problems, und extremale Graphentheorie mit dem Turán-Satz
und dem Szemerédi-Regularitätslemma. Ein zentrales Thema ist das
\emph{Hadwiger-Nelson-Problem}: Wie viele Farben werden benötigt, um die Ebene
so zu färben, dass keine zwei Punkte im Abstand 1 dieselbe Farbe erhalten?
Wir präsentieren die klassischen Schranken $4 \le \chr(\R^2) \le 7$, den
Moser-Spindle und den Golomb-Graphen, die $\chr(\R^2) \ge 4$ beweisen, und
das wegweisende Ergebnis von de Grey aus dem Jahr 2018, das $\chr(\R^2) \ge 5$
etabliert. Der genaue Wert von $\chr(\R^2)$ ist ein offenes Problem. Alle
Beweise werden vollständig präsentiert; Vermutungen werden klar von bewiesenen
Sätzen unterschieden.
\end{abstract}

\maketitle

\tableofcontents

%% =============================================================
\section{Einleitung}
%% =============================================================

Die Ursprünge der Graphentheorie liegen in einem berühmten Rätsel über die
Stadt Königsberg (heute Kaliningrad). Der Fluss Pregel teilte die Stadt in vier
Regionen auf, die durch sieben Brücken verbunden waren. Die Frage war, ob es
möglich sei, alle sieben Brücken genau einmal zu überqueren und zum Ausgangs\-punkt
zurückzukehren. Im Jahr 1736 bewies Leonhard Euler, dass kein solcher Weg
existiert, und begründete damit die Graphentheorie als eigenständige Disziplin
\cite{euler1736}.

Eulers Einsicht bestand darin, die Geographie vollständig zu abstrahieren:
die vier Landmassen werden durch \emph{Knoten} ersetzt und die sieben Brücken
durch \emph{Kanten}. Die Frage wird dadurch rein kombinatorischer Natur.
Euler beobachtete, dass ein geschlossener Weg, der jede Kante genau einmal
durchläuft (heute \emph{Eulerscher Kreis} genannt), genau dann existiert,
wenn jeder Knoten geraden Grad hat. Da Königsberg vier Knoten mit ungeradem
Grad hat, existiert kein solcher Kreis.

Von diesen bescheidenen Anfängen hat sich die Graphentheorie enorm ausgeweitet.
Heute umfasst sie:
\begin{itemize}[leftmargin=2em]
  \item \textbf{Strukturtheorie}: Bäume, Zusammenhang, Planarität, Minoren;
  \item \textbf{Färbetheorie}: Korrekte Färbungen, chromatische Polynome,
        Vierfarbensatz;
  \item \textbf{Spektraltheorie}: Eigenwerte von Adjazenz- und
        Laplace-Matrizen;
  \item \textbf{Extremale Theorie}: Turán-artige Probleme, Ramsey-Theorie;
  \item \textbf{Algorithmische Aspekte}: Matching, Netzwerkfluss, NP-Vollständigkeit;
  \item \textbf{Geometrische Graphentheorie}: Einheitsdistanzgraphen, das
        Hadwiger-Nelson-Problem.
\end{itemize}

Diese Arbeit ist wie folgt gegliedert. Abschnitt~\ref{sec:grundbegriffe}
führt die grundlegenden Definitionen und das Handschlag-Lemma ein.
Abschnitt~\ref{sec:baeume} behandelt Bäume und Spannbäume, einschließlich
der Cayley-Formel. Abschnitt~\ref{sec:zusammenhang} deckt Zusammenhang,
Menger-Satz und Max-Flow-Min-Cut ab. Abschnitt~\ref{sec:faerbung} entwickelt
die Graphenfärbetheorie und den Vierfarbensatz. Abschnitt~\ref{sec:planar}
untersucht planare Graphen, die Euler-Formel und den Kuratowski-Satz.
Abschnitt~\ref{sec:hadwiger} ist dem Hadwiger-Nelson-Problem gewidmet.
Abschnitt~\ref{sec:spektral} führt die spektrale Graphentheorie ein.
Abschnitt~\ref{sec:euler_hamilton} behandelt Eulersche und Hamiltonische
Graphen. Abschnitt~\ref{sec:extremal} deckt extremale Graphentheorie ab.
Abschnitt~\ref{sec:literatur} enthält die Literaturangaben.

\subsection*{Konventionen}

In dieser Arbeit bezeichnet $G = (V, E)$ einen einfachen ungerichteten Graphen
mit Knotenmenge $V$ und Kantenmenge $E \subseteq \binom{V}{2}$. Wir schreiben
$n = |V|$ und $m = |E|$. Ein \emph{einfacher} Graph hat keine Schleifen und
keine Mehrfachkanten. Wir bezeichnen den vollständigen Graphen auf $n$ Knoten
mit $K_n$, den vollständig bipartiten Graphen mit $K_{p,q}$, den Weg auf $n$
Knoten mit $P_n$ und den Kreis auf $n$ Knoten mit $C_n$.

%% =============================================================
\section{Grundbegriffe}
\label{sec:grundbegriffe}
%% =============================================================

\begin{definition}[Graph]
Ein \emph{einfacher ungerichteter Graph} ist ein Paar $G = (V, E)$, wobei
$V$ eine endliche nichtleere Menge von \emph{Knoten} ist und $E \subseteq
\binom{V}{2}$ eine Menge von \emph{Kanten}, von denen jede eine zweielementige
Teilmenge von $V$ ist. Ein \emph{gerichteter Graph} (Digraph) $\vec{G} = (V, A)$
hat eine Menge $A \subseteq V \times V$ von \emph{Bögen} (gerichteten Kanten),
wobei im Allgemeinen $(u, v) \ne (v, u)$ gilt.
\end{definition}

\begin{definition}[Grad]
Der \emph{Grad} $\grad(v)$ eines Knotens $v \in V$ ist die Anzahl der zu $v$
inzidenten Kanten:
\[
  \grad(v) = |\{u \in V : \{u,v\} \in E\}|.
\]
Der \emph{Minimalgrad} ist $\delta(G) = \min_{v \in V} \grad(v)$ und der
\emph{Maximalgrad} ist $\Delta(G) = \max_{v \in V} \grad(v)$.
\end{definition}

\begin{definition}[Adjazenzmatrix]
Die \emph{Adjazenzmatrix} von $G$ ist die $n \times n$-Matrix $A = (a_{ij})$,
wobei $a_{ij} = 1$ wenn $\{i,j\} \in E$ und $a_{ij} = 0$ sonst. Für einfache
Graphen ist $A$ symmetrisch mit Nulldiagonale.
\end{definition}

\begin{theorem}[Handschlag-Lemma]
\label{satz:handschlag}
Für jeden Graphen $G = (V, E)$ gilt:
\[
  \sum_{v \in V} \grad(v) = 2|E|.
\]
Insbesondere ist die Anzahl der Knoten mit ungeradem Grad stets gerade.
\end{theorem}

\begin{proof}
Jede Kante $\{u, v\} \in E$ trägt genau $1$ zu $\grad(u)$ und genau $1$ zu
$\grad(v)$ bei. Sie trägt daher genau $2$ zur Summe $\sum_{v \in V} \grad(v)$
bei. Summieren über alle $|E|$ Kanten ergibt $\sum_{v \in V} \grad(v) = 2|E|$.

Da $2|E|$ gerade ist und $\sum_{v \in V} \grad(v)$ eine Summe von ganzen Zahlen
gleich einer geraden Zahl ist, muss die Anzahl der ungeraden Summanden gerade
sein.
\end{proof}

\begin{definition}[Teilgraph, Isomorphismus]
Ein Graph $H = (W, F)$ ist ein \emph{Teilgraph} von $G = (V, E)$, wenn
$W \subseteq V$ und $F \subseteq E$. Er ist ein \emph{induzierter Teilgraph}
auf $S \subseteq V$, wenn $F = E \cap \binom{S}{2}$. Zwei Graphen $G$ und $H$
sind \emph{isomorph}, wenn es eine Bijektion $\phi: V(G) \to V(H)$ gibt,
sodass $\{u,v\} \in E(G)$ genau dann wenn $\{\phi(u), \phi(v)\} \in E(H)$.
\end{definition}

\begin{example}
Der vollständige Graph $K_n$ hat $n$ Knoten, $\binom{n}{2}$ Kanten und
jeder Knoten hat Grad $n-1$. Das Handschlag-Lemma ergibt
$n(n-1) = 2\binom{n}{2}$, was sich bestätigt.
\end{example}

\begin{definition}[Weg, Pfad, Kreis]
Ein \emph{Weg} in $G$ ist eine Folge $v_0, e_1, v_1, e_2, \ldots, e_k, v_k$,
wobei jedes $e_i = \{v_{i-1}, v_i\} \in E$. Ein \emph{Pfad} ist ein Weg ohne
wiederholte Knoten. Ein \emph{Kreis} ist ein geschlossener Weg ($v_0 = v_k$)
ohne weitere wiederholte Knoten, der Länge $k \ge 3$.
\end{definition}

%% =============================================================
\section{Bäume und Spannbäume}
\label{sec:baeume}
%% =============================================================

\begin{definition}[Baum, Wald, Spannbaum]
Ein \emph{Wald} ist ein kreisfreier Graph. Ein \emph{Baum} ist ein
zusammenhängender kreisfreier Graph. Ein \emph{Spannbaum} eines zusammen\-hängenden
Graphen $G$ ist ein Teilgraph $T \subseteq G$, der ein Baum ist und alle
Knoten von $G$ enthält.
\end{definition}

\begin{theorem}[Charakterisierung von Bäumen]
\label{satz:baum-char}
Für einen Graphen $G = (V, E)$ mit $|V| = n \ge 1$ sind folgende Aussagen
äquivalent:
\begin{enumerate}[label=\upshape(\roman*)]
  \item $G$ ist ein Baum (zusammenhängend und kreisfrei);
  \item $G$ ist zusammenhängend und $|E| = n - 1$;
  \item $G$ ist kreisfrei und $|E| = n - 1$;
  \item Je zwei verschiedene Knoten von $G$ sind durch genau einen Pfad
        verbunden;
  \item $G$ ist zusammenhängend, aber das Entfernen jeder Kante macht ihn
        nicht mehr zusammenhängend.
\end{enumerate}
\end{theorem}

\begin{proof}
Wir beweisen (i) $\Rightarrow$ (ii) $\Rightarrow$ (iii) $\Rightarrow$ (i)
und skizzieren die Äquivalenz mit (iv) und (v).

\medskip
\noindent\textbf{(i) $\Rightarrow$ (ii).}
Durch Induktion über $n$. Für $n = 1$: der einzige Baum auf einem Knoten
hat $0 = 1-1$ Kanten. Für $n \ge 2$: jeder Baum hat ein Blatt (einen Knoten
vom Grad 1). Sei $v$ ein Blatt von $T$. Dann ist $T - v$ ein zusammen\-hängender
kreisfreier Graph auf $n - 1$ Knoten, also nach Induktion
$|E(T-v)| = n - 2$. Da $T$ eine Kante mehr als $T - v$ hat (die zu $v$
inzidente Kante), erhalten wir $|E(T)| = n - 1$.

\medskip
\noindent\textbf{(ii) $\Rightarrow$ (iii).}
Angenommen, $G$ ist zusammenhängend mit $n - 1$ Kanten. Wir zeigen, dass $G$
kreisfrei ist. Ein zusammenhängender Graph mit $k$ Knoten hat mindestens $k - 1$
Kanten (nur bei Bäumen Gleichheit). Enthielte $G$ einen Kreis $C$, so würde
das Entfernen einer Kante von $C$ den Graphen $G$ zusammenhängend lassen und
zu einem zusammenhängenden Graphen auf $n$ Knoten mit $n-2$ Kanten führen.
Aber ein zusammenhängender Graph auf $n$ Knoten braucht mindestens $n-1$
Kanten --- Widerspruch.

\medskip
\noindent\textbf{(iii) $\Rightarrow$ (i).}
Angenommen, $G$ ist kreisfrei mit $n - 1$ Kanten. Ein Wald auf $n$ Knoten
mit $k$ zusammenhängenden Komponenten hat genau $n - k$ Kanten. Da $|E| = n - 1$,
erhalten wir $k = 1$, also ist $G$ zusammenhängend, mithin ein Baum.

\medskip
\noindent Die Äquivalenzen mit (iv) und (v) folgen, da eindeutige Pfade
kreisfreie zusammenhängende Graphen charakterisieren, und die Minimalität
des Zusammenhangs Bäume kennzeichnet.
\end{proof}

\begin{theorem}[Cayley-Formel]
\label{satz:cayley}
Die Anzahl der verschiedenen markierten Spannbäume von $K_n$ beträgt $n^{n-2}$.
\end{theorem}

\begin{proof}[Beweis (Prüfer-Folgen)]
Wir etablieren eine Bijektion zwischen markierten Spannbäumen auf
$[n] = \{1, \ldots, n\}$ und Folgen $(a_1, \ldots, a_{n-2}) \in [n]^{n-2}$.

\textbf{Kodierung.} Gegeben ein Baum $T$: wir entfernen iterativ das Blatt
mit dem kleinsten Label. In Schritt $i$ sei $\ell_i$ das Blatt mit kleinstem
Label, setze $a_i$ gleich dem eindeutigen Nachbarn von $\ell_i$ und entferne
$\ell_i$. Nach $n - 2$ Schritten verbleiben zwei Knoten; wir erhalten die
Folge $(a_1, \ldots, a_{n-2})$.

\textbf{Dekodierung.} Gegeben eine Folge $(a_1, \ldots, a_{n-2})$: in Schritt
$i$ ist das aktuelle Blatt das kleinste Element von $[n]$, das nicht in
$\{a_i, a_{i+1}, \ldots, a_{n-2}\}$ und noch nicht entfernt wurde; füge eine
Kante von diesem Blatt zu $a_i$ hinzu.

Man verifiziert, dass Kodierung und Dekodierung inverse Bijektionen zwischen
der Menge der markierten Spannbäume von $K_n$ und $[n]^{n-2}$ sind. Da
$|[n]^{n-2}| = n^{n-2}$, folgt die Formel.
\end{proof}

\begin{remark}
Für $n = 4$: $4^2 = 16$ markierte Spannbäume von $K_4$. Für $n = 2$: $2^0 = 1$
(die einzelne Kante). Diese stimmen mit direkter Aufzählung überein.
\end{remark}

%% =============================================================
\section{Zusammenhang}
\label{sec:zusammenhang}
%% =============================================================

\begin{definition}[Zusammenhang, Schnitt]
Ein Graph $G$ ist \emph{$k$-zusammenhängend} ($k$-knotenzusammenhängend),
wenn $|V| > k$ und das Entfernen beliebiger $k-1$ Knoten $G$ zusammenhängend
lässt. Der \emph{Knotenzusammenhang} $\kappa(G)$ ist das maximale solche $k$.
Ein \emph{$s$-$t$-Knotenschnitt} ist eine Menge $S \subset V \setminus \{s,t\}$,
deren Entfernung $s$ von $t$ trennt.
\end{definition}

\begin{theorem}[Menger-Satz]
\label{satz:menger}
Sei $G$ ein Graph und $s, t \in V$ mit $\{s,t\} \notin E$. Die maximale Anzahl
paarweise intern knotendisjunkter $s$-$t$-Pfade ist gleich der minimalen Größe
eines $s$-$t$-Knotenschnitts.
\end{theorem}

\begin{proof}[Beweisskizze]
Sei $f^*$ die maximale Anzahl knotendisjunkter $s$-$t$-Pfade und $c^*$ die
minimale Knotenschnittgröße. Offensichtlich $f^* \le c^*$ (jeder Schnitt
muss jeden Pfad ``blockieren''). Für die andere Richtung verwenden wir ein
Flussargument.

Konstruiere ein gerichtetes Netzwerk $N$ aus $G$: ersetze jeden Knoten
$v \ne s, t$ durch zwei Knoten $v^{\text{ein}}$ und $v^{\text{aus}}$, die
durch einen Bogen der Kapazität $1$ verbunden sind; ersetze jede Kante
$\{u,v\}$ durch Bögen $(u^{\text{aus}}, v^{\text{ein}})$ und
$(v^{\text{aus}}, u^{\text{ein}})$, jeweils mit Kapazität $\infty$. Ein
maximaler ganzzahliger $s$-$t$-Fluss in $N$ entspricht knotendisjunkten
Pfaden in $G$, und ein minimaler $s$-$t$-Schnitt in $N$ entspricht einem
minimalen Knotenschnitt in $G$. Nach dem Max-Flow-Min-Cut-Satz
(Satz~\ref{satz:maxflow}) sind diese Werte gleich.
\end{proof}

\begin{theorem}[Max-Flow-Min-Cut-Satz, Ford-Fulkerson]
\label{satz:maxflow}
In einem gerichteten Netzwerk $N = (V, A, c)$ mit Quelle $s$ und Senke $t$
und nichtnegativen ganzzahligen Kapazitäten $c: A \to \Z_{\ge 0}$ ist der
Wert eines maximalen $s$-$t$-Flusses gleich der Kapazität eines minimalen
$s$-$t$-Schnitts.
\end{theorem}

\begin{proof}
Sei $f$ ein beliebiger $s$-$t$-Fluss und $(S, T)$ ein beliebiger $s$-$t$-Schnitt
(mit $s \in S$, $t \in T$, $S \cup T = V$, $S \cap T = \emptyset$). Nach der
Flusserhaltung gilt:
$|f| = \sum_{(u,v): u \in S, v \in T} f(u,v) - \sum_{(u,v): u \in T, v \in S} f(u,v)
\le \sum_{(u,v): u \in S, v \in T} c(u,v) = c(S,T)$.
Also $\max |f| \le \min c(S,T)$.

Für die Gleichheit führen wir den Ford-Fulkerson-Algorithmus aus: suche
wiederholt einen \emph{augmentierenden Pfad} von $s$ nach $t$ im Residualgraphen
und erhöhe den Fluss entlang ihm. Dies terminiert für ganzzahlige Kapazitäten.
Wenn kein augmentierender Pfad mehr existiert, definiere $S$ als die Menge
der von $s$ im Residualgraphen erreichbaren Knoten. Dann ist $(S, V \setminus S)$
ein Schnitt der Kapazität $|f|$. Also $\max |f| = \min c(S,T)$.
\end{proof}

%% =============================================================
\section{Graphenfärbung}
\label{sec:faerbung}
%% =============================================================

\begin{definition}[Korrekte Färbung, Chromatische Zahl]
Eine \emph{korrekte $k$-Färbung} von $G$ ist eine Funktion $c: V \to \{1, \ldots, k\}$,
sodass $c(u) \ne c(v)$ wann immer $\{u,v\} \in E$. Die \emph{chromatische Zahl}
$\chr(G)$ ist das kleinste $k$, für das eine korrekte $k$-Färbung existiert.
\end{definition}

\begin{example}
$\chr(K_n) = n$, $\chr(C_{2k}) = 2$, $\chr(C_{2k+1}) = 3$, $\chr(P_n) = 2$
für $n \ge 2$.
\end{example}

\begin{proposition}[Greedy-Schranke]
\label{prop:greedy}
Für jeden Graphen $G$ gilt $\chr(G) \le \Delta(G) + 1$.
\end{proposition}

\begin{proof}
Ordne die Knoten beliebig als $v_1, \ldots, v_n$. Weise jedem Knoten $v_i$
die kleinste Farbe aus $\{1, 2, \ldots\}$ zu, die von keinem seiner bereits
gefärbten Nachbarn verwendet wird. Da $v_i$ höchstens $\Delta(G)$ Nachbarn
hat, sind höchstens $\Delta(G)$ Farben ausgeschlossen, sodass stets eine
Farbe aus $\{1, \ldots, \Delta(G)+1\}$ verfügbar ist.
\end{proof}

\begin{theorem}[Brook-Satz]
\label{satz:brook}
Ist $G$ ein zusammenhängender Graph, der weder ein vollständiger Graph noch
ein ungerader Kreis ist, so gilt $\chr(G) \le \Delta(G)$.
\end{theorem}

\begin{proof}[Beweisskizze]
Wir führen Induktion über $|V|$ durch. Für $\Delta(G) = 2$ muss $G$ ein Pfad
oder ein gerader Kreis sein (ungerade Kreise und Dreiecke sind ausgeschlossen),
beide 2-färbbar.

Für $\Delta(G) \ge 3$: Ist $G$ nicht $\Delta$-regulär, entfernen wir einen
Knoten $v$ mit $\grad(v) < \Delta(G)$. Der verbleibende Graph $G - v$ kann
nach Induktion $\Delta$-gefärbt werden, und da $v$ weniger als $\Delta$ Nachbarn
hat, ist für $v$ noch eine Farbe frei.

Ist $G$ $\Delta$-regulär mit einem Trennknoten, zerlegen wir den Graphen
und wenden Induktion an. Ist $G$ $\Delta$-regulär und 2-zusammenhängend,
ermöglicht eine sorgfältige Knotenumordnung (mittels eines Pfades und
Rückweges) dem Greedy-Algorithmus, mit $\Delta$ Farben auszukommen.
\end{proof}

\begin{theorem}[Vierfarbensatz]
\label{satz:vierfarbensatz}
Jeder planare Graph ist $4$-färbbar: $\chr(G) \le 4$ für alle planaren
Graphen $G$.
\end{theorem}

\begin{proof}[Historische Anmerkung und Beweisskizze]
Der Vierfarbensatz hat eine reiche Geschichte. Francis Guthrie vermutete
ihn 1852. Alfred Kempe veröffentlichte 1879 einen vermeintlichen Beweis,
den Percy Heawood 1890 widerlegte (er fand eine Lücke); Kempes Argument
lieferte jedoch noch immer den Fünffarbensatz
(Satz~\ref{satz:fuenffarbensatz}).

Der Satz wurde schließlich \textbf{von Appel und Haken 1976 bewiesen}
\cite{appelhaken1977}, mithilfe einer unvermeidbaren Menge von 1936
reduzierbaren Konfigurationen, die durch Computer verifiziert wurden. Dies
war das erste bedeutende Theorem, dessen Beweis wesentlich auf
computergestützte Hilfe angewiesen war.

Ein zweiter, saubererer Beweis wurde \textbf{1997 von Robertson, Sanders,
Seymour und Thomas} \cite{robertson1997} gegeben, der die unvermeidbare
Menge auf 633 Konfigurationen reduziert und die Computerverifikation
transparenter gestaltet.

\textbf{Beweissstrategie.} Das Argument verläuft in zwei Schritten:
\begin{enumerate}
  \item \textbf{Unvermeidbarkeit:} Zeige, dass jeder planare Graph mindestens
        eine Konfiguration aus einer endlichen Liste enthält (z.\,B.\
        Konfigurationen von Patches mit kleinen Flächen um einen Knoten).
  \item \textbf{Reduzierbarkeit:} Zeige, dass, wenn eine solche Konfiguration
        in einem minimalen Gegenbeispiel erscheint, ein kleinerer planarer
        Graph, der nicht 4-färbbar ist, abgeleitet werden kann ---
        im Widerspruch zur Minimalität.
\end{enumerate}
Die Kombination von Unvermeidbarkeit und Reduzierbarkeit zeigt, dass kein
minimales Gegenbeispiel existieren kann, also auch kein Gegenbeispiel
überhaupt.
\end{proof}

\begin{remark}
Kein elementarer Beweis (ohne Computerunterstützung) des Vierfarbensatzes
ist bekannt. Beide bekannten Beweise sind computergestützt.
\end{remark}

%% =============================================================
\section{Planare Graphen}
\label{sec:planar}
%% =============================================================

\begin{definition}[Planarer Graph, Einbettung, Fläche]
Ein Graph $G$ ist \emph{planar}, wenn er in der Ebene $\R^2$ so gezeichnet
werden kann, dass sich keine zwei Kanten außer an gemeinsamen Endpunkten
kreuzen. Eine solche Zeichnung ist eine \emph{planare Einbettung}. Eine
planare Einbettung teilt die Ebene in zusammenhängende Regionen, sogenannte
\emph{Flächen}, einschließlich genau einer unbeschränkten Fläche.
\end{definition}

\begin{theorem}[Euler-Formel]
\label{satz:eulerformel}
Sei $G$ ein zusammenhängender planarer Graph mit $V$ Knoten, $E$ Kanten und
$F$ Flächen in einer planaren Einbettung. Dann gilt:
\[
  V - E + F = 2.
\]
\end{theorem}

\begin{proof}
Per Induktion über die Anzahl der Kanten. Ist $G$ ein Baum, dann $E = V - 1$,
$F = 1$ (nur die unbeschränkte Fläche), und $V - (V-1) + 1 = 2$. ✓

Angenommen, $G$ hat einen Kreis, also existiert eine Kante $e$ auf einem
Kreis. Das Entfernen von $e$ vereinigt zwei Flächen zu einer, also
$F' = F - 1$, $E' = E - 1$, $V' = V$. Nach Induktion gilt $V' - E' + F' = 2$,
d.\,h.\ $V - (E-1) + (F-1) = 2$, was $V - E + F = 2$ ergibt. ✓

Ist $G$ zusammenhängend, hat aber eine Brücke $e$ (die auf keinem Kreis liegt),
so trennt das Entfernen von $e$ den Graphen: $V' = V$, $E' = E - 1$. Die zwei
resultierenden Komponenten $G_1, G_2$ erfüllen $V_1 - E_1 + F_1 = 2$ und
$V_2 - E_2 + F_2 = 2$. Die ursprüngliche Einbettung hat die geteilte Fläche
einmal gezählt, also $F = F_1 + F_2 - 1$, $V = V_1 + V_2$, $E = E_1 + E_2 + 1$.
Addition der beiden Gleichungen und Anpassung ergibt $V - E + F = 2$.
\end{proof}

\begin{corollary}[Kantenschranke für planare Graphen]
\label{kor:planare-kanten}
Ist $G$ ein einfacher planarer Graph mit $n \ge 3$ Knoten und $m$ Kanten, so
gilt $m \le 3n - 6$. Ist $G$ zusätzlich dreiecksfrei, so gilt $m \le 2n - 4$.
\end{corollary}

\begin{proof}
Jede Fläche ist von mindestens 3 Kanten begrenzt. Da jede Kante genau
2 Flächen begrenzt, gilt $3F \le 2E$, also $F \le 2E/3$. Einsetzen in
die Euler-Formel: $2 = V - E + F \le n - E + 2E/3 = n - E/3$, was
$E \le 3n - 6$ ergibt. Für dreiecksfreie Graphen ist jede Fläche von
mindestens 4 Kanten begrenzt, was $4F \le 2E$ und damit $E \le 2n - 4$ liefert.
\end{proof}

\begin{corollary}
$K_5$ und $K_{3,3}$ sind nicht planar.
\end{corollary}

\begin{proof}
Für $K_5$: $n = 5$, $m = 10 > 3 \cdot 5 - 6 = 9$. Verstößt gegen
Korollar~\ref{kor:planare-kanten}.
Für $K_{3,3}$: $n = 6$, $m = 9$. Es ist bipartit, also dreiecksfrei.
Aber $9 > 2 \cdot 6 - 4 = 8$. Verstößt gegen die dreiecksfreie Schranke.
\end{proof}

\begin{theorem}[Kuratowski-Satz]
\label{satz:kuratowski}
Ein Graph $G$ ist planar genau dann, wenn er keine Unterteilung von $K_5$ oder
$K_{3,3}$ als Teilgraph enthält.
\end{theorem}

\begin{proof}[Beweisskizze]
Die Notwendigkeit (planar $\Rightarrow$ keine $K_5$- oder $K_{3,3}$-Unterteilung)
folgt aus Korollar~\ref{kor:planare-kanten}, angewandt auf Unterteilungen.

Für die Hinlänglichkeit geht man durch Widerspruch vor: Man nehme an, $G$ sei
ein minimaler nicht-planarer Graph. Mithilfe der Eigenschaften minimaler
nicht-planarer Graphen (sie sind 3-zusammenhängend; jede Kante liegt auf einem
Kreis) zeigt man, dass $G$ eine $K_5$- oder $K_{3,3}$-Unterteilung enthalten
muss. Der vollständige Beweis findet sich in \cite{diestel2017}.
\end{proof}

\begin{theorem}[Fünffarbensatz]
\label{satz:fuenffarbensatz}
Jeder planare Graph ist $5$-färbbar.
\end{theorem}

\begin{proof}
Durch Induktion über $n = |V|$. Die Basisfälle sind trivial. Für den
Induktionsschritt sei $G$ ein planarer Graph mit $n$ Knoten. Nach
Korollar~\ref{kor:planare-kanten} hat $G$ einen Knoten $v$ mit $\grad(v) \le 5$.

\textbf{Fall 1: $\grad(v) \le 4$.} Entferne $v$, wende Induktion an, um eine
5-Färbung von $G - v$ zu erhalten. Da $v$ höchstens 4 Nachbarn hat, ist
mindestens eine der 5 Farben für $v$ frei.

\textbf{Fall 2: $\grad(v) = 5$.} Seien die Nachbarn $v_1, \ldots, v_5$ in
zyklischer Reihenfolge um $v$ in der Einbettung. Haben zwei nicht-adjazente
Nachbarn in der $(G-v)$-Färbung dieselbe Farbe, wird $v$ mit der freigewordenen
Farbe gefärbt.

Andernfalls verwenden $v_1, \ldots, v_5$ alle 5 verschiedenen Farben. Betrachte
den $v_1$-$v_3$-Teilgraphen, der durch die Farben 1 und 3 induziert wird.
Liegen $v_1$ und $v_3$ in verschiedenen Zusammenhangskomponenten dieses
2-gefärbten Teilgraphen, tausche die Farben in $v_1$s Komponente. Jetzt hat
$v_1$ die Farbe 3, was Farbe 1 für $v$ freigibt.

Liegen $v_1$ und $v_3$ in derselben Komponente, gibt es einen Pfad $P$ von
$v_1$ nach $v_3$, der nur die Farben 1 und 3 verwendet. Zusammen mit den
Kanten $vv_1$ und $vv_3$ trennt $P$ den Knoten $v_2$ von $v_4, v_5$. Dann
liegen $v_2$ und $v_4$ in verschiedenen Komponenten des 2-gefärbten Teilgraphen
mit den Farben 2 und 4, sodass wir tauschen können, um Farbe 2 für $v$
freizugeben.
\end{proof}

%% =============================================================
\section{Das Hadwiger-Nelson-Problem}
\label{sec:hadwiger}
%% =============================================================

\begin{definition}[Einheitsdistanzgraph, Chromatische Zahl der Ebene]
Der \emph{Einheitsdistanzgraph der Ebene} $G_{\R^2, 1}$ hat die Knotenmenge
$\R^2$ und eine Kante zwischen je zwei Punkten mit euklidischem Abstand $1$.
Die \emph{chromatische Zahl der Ebene} ist:
\[
  \chr(\R^2) = \chr(G_{\R^2, 1})
  = \min\bigl\{k : \exists\; k\text{-Färbung von } \R^2 \text{ ohne zwei
    Einheitsdistanz-Punkte gleicher Farbe}\bigr\}.
\]
\end{definition}

Das \emph{Hadwiger-Nelson-Problem}, unabhängig von Hugo Hadwiger, Edward Nelson
und John Isbell um 1950 gestellt, fragt nach dem genauen Wert von $\chr(\R^2)$.

\begin{theorem}[Klassische Schranken]
\label{satz:hn-klassisch}
$4 \le \chr(\R^2) \le 7.$
\end{theorem}

\begin{proof}
\textbf{Obere Schranke $\chr(\R^2) \le 7$.}
Kachele die Ebene mit regulären Sechsecken der Seitenlänge $s = 0{,}48 <
1/\sqrt{3} \approx 0{,}577$. Färbe die Sechsecke mit 7 Farben, indem man das
periodische Muster $(i,j) \mapsto (i + 2j) \bmod 7$ in den hexagonalen
Gitterkoordinaten verwendet.

Zwei Punkte im selben Sechseck haben Abstand höchstens $2s = 0{,}96 < 1$.
Zwei Punkte in gleichfarbigen Sechsecken haben Mittelpunktabstand mindestens
$\sqrt{7} \cdot s\sqrt{3} = \sqrt{21} \cdot 0{,}48 \approx 2{,}20 > 1$.
Also haben gleichfarbige Punkte stets Abstand $\ne 1$, was $\chr(\R^2) \le 7$
beweist.

\textbf{Untere Schranke $\chr(\R^2) \ge 4$.}
Der \emph{Moser-Spindle} (Moser \& Moser, 1961) ist ein Einheitsdistanzgraph
mit 7 Knoten und 11 Kanten. Sei er mit Knoten $A, B, C, D, E, F, G$ in $\R^2$
eingebettet, wobei alle Kanten die Länge 1 haben. Der Graph hat die Eigenschaft,
dass $A$ und $G$ durch einen Pfad $A - B - C - D$ und einen weiteren Pfad
$A - E - F - G$ verbunden sind, wobei $D$ und $G$ durch eine Kante verbunden
sind. Eine Fallunterscheidung zeigt, dass jede 3-Färbung $A$ und $G$ zur
gleichen Farbe zwingt, obwohl $\{A, G\} \in E$ --- Widerspruch. Also
$\chr(\text{Moser-Spindle}) = 4 \le \chr(\R^2)$.
\end{proof}

\subsection{Der Moser-Spindle}

Der Moser-Spindle wird wie folgt konstruiert. Sei:
\begin{align*}
  A &= (0, 0), \quad B = (1, 0), \quad C = (1/2,\, \sqrt{3}/2),\\
  D &= (3/2,\, \sqrt{3}/2), \quad E = (\cos\theta,\, \sin\theta)
      \text{ mit } \cos\theta = 5/6,\\
  F &= \bigl((5-\sqrt{33})/12,\, (\sqrt{11}+5\sqrt{3})/12\bigr),\\
  G &= \bigl((15-\sqrt{33})/12,\, (3\sqrt{11}+5\sqrt{3})/12\bigr).
\end{align*}
Die Kanten sind: $AB, AC, BC$ (Dreieck $ABC$), $BD, CD$ ($D$ adjazent zu
$B$ und $C$), $AE, AF, EF$ (Dreieck $AEF$), $EG, FG$ ($G$ adjazent zu
$E$ und $F$) und $DG$.

Man verifiziert numerisch, dass alle 11 Kanten die Länge genau 1 haben.

\begin{lemma}[Chromatische Zahl des Moser-Spindle]
$\chr(\text{Moser-Spindle}) = 4$.
\end{lemma}

\begin{proof}
Der Graph ist nicht 3-färbbar. Angenommen, es existiere eine gültige
3-Färbung mit Farben $\{1, 2, 3\}$. Da $ABC$ ein Dreieck ist, erhalten
$A$, $B$, $C$ drei verschiedene Farben; sei $c(A) = 1$, $c(B) = 2$,
$c(C) = 3$. Da $BD$ und $CD$ Kanten sind, gilt $c(D) \ne 2$ und
$c(D) \ne 3$, also $c(D) = 1 = c(A)$. Da $AEF$ ein Dreieck ist, gilt
$c(E), c(F) \in \{2, 3\}$ mit $c(E) \ne c(F)$. Da $EG$ und $FG$ Kanten
sind, gilt $c(G) \ne c(E)$ und $c(G) \ne c(F)$, also muss $c(G)$ von
beiden Farben in $\{2,3\}$ verschieden sein, was $c(G) = 1$ erzwingt.
Aber $DG$ ist eine Kante und $c(D) = c(G) = 1$: Widerspruch. Damit gilt
$\chr(\text{Moser-Spindle}) \ge 4$. Der Graph ist 4-färbbar (z.\,B.\ weise
die Farben $1,2,3,1,2,3,4$ den Knoten $A,\ldots,G$ zu, die Kanten
respektierend), also $\chr = 4$.
\end{proof}

\subsection{De Greys Ergebnis von 2018}

\begin{theorem}[De Grey, 2018]
\label{satz:degrey}
$\chr(\R^2) \ge 5$.
\end{theorem}

\begin{proof}[Beweisskizze]
Aubrey de Grey konstruierte einen endlichen Einheitsdistanzgraphen $H$ mit
$\chr(H) = 5$ \cite{degrey2018}. Die Konstruktion begann mit dem Moser-Spindle
und wendete eine Folge von Operationen (Rotation, Vereinigung, Teilgraphenbildung)
an, um Graphen zu erhalten, die immer mehr Farben erfordern. Nach umfangreicher
Computersuche fand de Grey einen Graphen mit $\mathbf{1581}$ Knoten und
$\mathbf{12.066}$ Kanten, der nicht 4-färbbar ist. Dieser wurde von mehreren
unabhängigen Teams verifiziert und 2019 in \emph{Geombinatorics}
peer-reviewed publiziert.

Da $H$ ein Einheitsdistanzgraph in $\R^2$ ist, induziert jede Färbung von
$\R^2$ eine Färbung von $H$. Wäre $\chr(\R^2) \le 4$, so wäre $H$
4-färbbar --- Widerspruch. Also $\chr(\R^2) \ge 5$.
\end{proof}

\begin{conjecture}[Genauer Wert von $\chr(\R^2)$]
\label{verm:hn-genau}
Die chromatische Zahl der Ebene erfüllt $\chr(\R^2) \in \{5, 6, 7\}$.
Der genaue Wert ist unbekannt. De Greys Ergebnis schließt $\chr(\R^2) \le 4$
aus, und das hexagonale Färbungsschema zeigt $\chr(\R^2) \le 7$, aber ob
$\chr(\R^2) = 5$, $6$ oder $7$ ist, bleibt ein offenes Problem.
\end{conjecture}

\begin{remark}
Das Problem ist bemerkenswert, da die Antwort nicht vom Auswahlaxiom abhängt:
jede Färbung, die $\chr(\R^2) = k$ bezeugt, kann messbar gewählt werden, und
das Problem macht Sinn in einem rein kombinatorischen Rahmen. Das
Polymath-Projekt (2018) reduzierte anschließend die Knotenzahl eines
5-chromatischen Einheitsdistanzgraphen auf 553 Knoten.
\end{remark}

%% =============================================================
\section{Spektrale Graphentheorie}
\label{sec:spektral}
%% =============================================================

\begin{definition}[Laplace-Matrix]
Die \emph{Gradmatrix} von $G$ ist $D = \mathrm{diag}(\grad(v_1), \ldots,
\grad(v_n))$. Die \emph{Laplace-Matrix} ist
\[
  L = D - A,
\]
wobei $A$ die Adjazenzmatrix ist. Explizit: $L_{ij} = -1$ wenn
$\{i,j\} \in E$, $L_{ii} = \grad(v_i)$ und $L_{ij} = 0$ sonst.
\end{definition}

\begin{proposition}[Eigenschaften der Laplace-Matrix]
\label{prop:laplace}
Sei $G$ ein Graph mit $n$ Knoten und Laplace-Matrix $L$. Dann gilt:
\begin{enumerate}[label=\upshape(\roman*)]
  \item $L$ ist symmetrisch und positiv semidefinit;
  \item $L \mathbf{1} = \mathbf{0}$, also ist $0$ stets ein Eigenwert mit
        Eigenvektor $\mathbf{1} = (1,\ldots,1)^\top$;
  \item Die Anzahl der Null-Eigenwerte von $L$ ist gleich der Anzahl der
        Zusammenhangskomponenten von $G$;
  \item Für jeden Vektor $x \in \R^n$:
        $x^\top L x = \sum_{\{i,j\} \in E} (x_i - x_j)^2 \ge 0$.
\end{enumerate}
\end{proposition}

\begin{proof}
(iv) ist eine direkte Berechnung:
$x^\top L x = x^\top D x - x^\top A x = \sum_i \grad(v_i) x_i^2
- 2\sum_{\{i,j\} \in E} x_i x_j = \sum_{\{i,j\} \in E}(x_i^2 - 2x_ix_j + x_j^2)
= \sum_{\{i,j\} \in E}(x_i - x_j)^2 \ge 0$.
Dies beweist (i). (ii) folgt aus $\sum_j L_{ij} = \grad(v_i) - \sum_{j: \{i,j\} \in E} 1 = 0$.
(iii) folgt, da der Kern von $L$ den lokal konstanten Funktionen auf $G$
entspricht, also den Funktionen, die auf jeder Zusammenhangskomponente
konstant sind.
\end{proof}

\begin{definition}[Algebraische Konnektivität, Fiedler-Wert]
Die Eigenwerte von $L$ seien $0 = \lambda_1 \le \lambda_2 \le \cdots \le \lambda_n$.
Für einen zusammenhängenden Graphen gilt $\lambda_2 > 0$; dieser Wert heißt
die \emph{algebraische Konnektivität} oder der \emph{Fiedler-Wert} von $G$,
bezeichnet mit $a(G)$. Der zugehörige Eigenvektor ist der \emph{Fiedler-Vektor}.
\end{definition}

\begin{definition}[Kantenexpansion, Cheeger-Konstante]
Die \emph{Kantenexpansion} oder \emph{Cheeger-Konstante} von $G$ ist
\[
  h(G) = \min_{\emptyset \ne S \subsetneq V,\, |S| \le n/2}
         \frac{|E(S, V \setminus S)|}{|S|},
\]
wobei $E(S, V\setminus S)$ die Menge der Kanten mit einem Endpunkt in $S$
und dem anderen in $V \setminus S$ ist.
\end{definition}

\begin{theorem}[Cheeger-Ungleichung]
\label{satz:cheeger}
Für jeden $d$-regulären Graphen $G$ gilt:
\[
  \frac{h(G)^2}{2d} \le \lambda_2(L) \le 2h(G).
\]
\end{theorem}

\begin{proof}[Beweisskizze]
Die obere Schranke $\lambda_2(L) \le 2h(G)$ folgt aus der Rayleigh-Quotienten-Charakterisierung
von $\lambda_2$: für jeden zu $\mathbf{1}$ orthogonalen Vektor $x$ gilt
$\lambda_2 \le x^\top L x / x^\top x$. Wähle $x$ als die charakteristische
Funktion der optimalen Menge $S$, bei 0 zentriert; die resultierende Schranke
ergibt $\lambda_2 \le 2h(G)$.

Die untere Schranke $h(G) \le \sqrt{2d \lambda_2(L)}$ verwendet die
Co-Area-Formel (diskretes Federer-Fleming-Theorem), um die Kantenbegrenzung
von Niveaumengen des Fiedler-Vektors mit der algebraischen Konnektivität
in Beziehung zu setzen.
\end{proof}

\begin{remark}
Die Cheeger-Ungleichung ist das diskrete Analogon der isoperimetrischen
Ungleichung in der Riemannschen Geometrie. Sie verbindet kombinatorische
Expansion (nützlich für schnelles Mischen von Zufallsläufen) mit dem
Spektralgap.
\end{remark}

%% =============================================================
\section{Eulersche und Hamiltonische Graphen}
\label{sec:euler_hamilton}
%% =============================================================

\begin{definition}[Eulerscher Kreis, Hamiltonischer Kreis]
Ein \emph{Eulerscher Weg} ist ein Weg, der jede Kante genau einmal durchläuft.
Ein \emph{Eulerscher Kreis} ist ein geschlossener Eulerscher Weg. Ein Graph,
der einen Eulerschen Kreis enthält, heißt \emph{Eulersch}. Ein
\emph{Hamiltonischer Kreis} ist ein Kreis, der jeden Knoten genau einmal
besucht. Ein Graph, der einen Hamiltonischen Kreis enthält, heißt
\emph{Hamiltonisch}.
\end{definition}

\begin{theorem}[Eulers Charakterisierung Eulerscher Graphen]
\label{satz:eulersch}
Ein zusammenhängender Graph $G$ ist Eulersch genau dann, wenn jeder Knoten
geraden Grad hat.
\end{theorem}

\begin{proof}
\textbf{Notwendigkeit.} Bei einem Eulerschen Kreis wird jedes Mal, wenn der
Weg einen Knoten $v$ besucht, eine Kante für die Ankunft und eine für die
Abfahrt verwendet, was $2$ zu $\grad(v)$ beiträgt. Damit ist $\grad(v)$ gerade
für alle $v$.

\textbf{Hinlänglichkeit.} Sei $G$ zusammenhängend mit allen Graden gerade.
Beginne einen Weg bei einem beliebigen Knoten $v_0$. Der Weg kann stets
fortgesetzt werden (da gerade Grade sicherstellen, dass eine Ankunft bei einem
Knoten die Existenz einer unbenutzten Abfahrtskante impliziert), bis er sich
bei $v_0$ schließt. Benutzt dieser geschlossene Weg nicht alle Kanten, gibt
es eine unbenutzte Kante $e = \{u, w\}$, die zu einem Knoten $u$ auf dem Weg
inzident ist (da $G$ zusammenhängend ist). Beginne bei $u$ einen neuen
geschlossenen Weg, der nur unbenutzte Kanten verwendet; füge diesen in den
ursprünglichen Weg bei $u$ ein. Wiederhole, bis alle Kanten benutzt sind.
\end{proof}

\begin{remark}
Für das Königsberger Brückenproblem: Die vier Landmassen haben ungerade Grade
(3, 3, 3, 5), also existiert kein Eulerscher Kreis. Dies bestätigt Eulers
Ergebnis von 1736.
\end{remark}

\begin{theorem}[Satz von Dirac]
\label{satz:dirac}
Ist $G$ ein einfacher Graph auf $n \ge 3$ Knoten mit $\delta(G) \ge n/2$,
so ist $G$ Hamiltonisch.
\end{theorem}

\begin{proof}
Sei $P = v_1, v_2, \ldots, v_k$ ein längster Pfad in $G$. Wir behaupten,
$k = n$ und der Pfad kann zu einem Hamiltonischen Kreis geschlossen werden.

Ist $k < n$, so hat $v_1$ mindestens $n/2$ Nachbarn, alle auf $P$ (da $P$
maximal ist). Ebenso für $v_k$. Durch ein Zählargument gibt es Indizes $i$,
sodass $v_k v_i \in E$ und $v_1 v_{i+1} \in E$. Damit kann man einen Kreis
$C = v_1 \cdots v_i v_k v_{k-1} \cdots v_{i+1} v_1$ bilden. Da $G$
zusammenhängend ist und $k < n$, ist ein Knoten $w \notin V(C)$ zu einem
Knoten in $C$ adjazent, was eine Verlängerung von $P$ ermöglicht --- im
Widerspruch zur Maximalität. Also $k = n$ und $G$ ist Hamiltonisch.
\end{proof}

\begin{theorem}[NP-Vollständigkeit des Hamiltonischen-Pfad-Problems]
\label{satz:ham-np}
Das Problem, zu entscheiden, ob ein gegebener Graph $G$ einen Hamiltonischen
Pfad (oder Kreis) enthält, ist NP-vollständig.
\end{theorem}

\begin{proof}
Das Problem liegt in NP: Ein Hamiltonischer Pfad kann, wenn er gegeben ist,
in polynomieller Zeit verifiziert werden (prüfe, dass jeder Knoten genau
einmal vorkommt und alle Kanten in $E$ liegen).

NP-Schwerheit wird durch Reduktion von 3-SAT gezeigt. Gegeben eine
3-KNF-Formel $\phi$ auf $n$ Variablen und $m$ Klauseln, konstruiere einen
Graphen $G_\phi$, sodass $G_\phi$ einen Hamiltonischen Kreis hat genau dann,
wenn $\phi$ erfüllbar ist. Die Standard-Reduktion (Karp, 1972 \cite{karp1972})
verwendet ein ``Variablen-Gadget'' für jede Variable (ein Paar von Pfaden
für wahr/falsch) und ein ``Klauseln-Gadget'' für jede Klausel (ein einzelner
Knoten, der mit den Variablen-Pfaden verbunden ist). Der resultierende Graph
hat $O(nm)$ Knoten und Kanten, konstruierbar in polynomieller Zeit. Die
Korrektheit der Reduktion sichert die Äquivalenz.
\end{proof}

\begin{remark}
Das \emph{Travelling-Salesman-Problem} (TSP) fragt nach dem
minimalen gewichteten Hamiltonischen Kreis in einem gewichteten vollständigen
Graphen. TSP ist NP-schwer und ein zentrales Problem der kombinatorischen
Optimierung. Für das euklidische TSP (Punkte in der Ebene, euklidische
Abstände) existiert eine $(1+\varepsilon)$-Approximation in polynomieller
Zeit für jedes $\varepsilon > 0$ (PTAS, Arora 1998), aber exakte Lösung
bleibt NP-schwer.
\end{remark}

%% =============================================================
\section{Extremale Graphentheorie}
\label{sec:extremal}
%% =============================================================

Die extremale Graphentheorie fragt: Wie viele Kanten kann ein Graph auf $n$
Knoten haben, wenn ein bestimmter Teilgraph verboten ist?

\begin{definition}[Turán-Graph]
Der \emph{Turán-Graph} $T(n, r)$ ist der vollständige $r$-partite Graph auf
$n$ Knoten, dessen Teile so gleich wie möglich sind (Größen
$\lfloor n/r \rfloor$ oder $\lceil n/r \rceil$). Seine Kantenzahl ist
\[
  t(n, r) = \left(1 - \frac{1}{r}\right)\frac{n^2}{2} + O(n).
\]
\end{definition}

\begin{theorem}[Turán-Satz]
\label{satz:turan}
Hat ein einfacher Graph $G$ auf $n$ Knoten mehr als $t(n, r)$ Kanten, so
enthält $G$ den vollständigen Graphen $K_{r+1}$ als Teilgraph. Umgekehrt ist
$T(n,r)$ der eindeutig dichteste $K_{r+1}$-freie Graph auf $n$ Knoten.
\end{theorem}

\begin{proof}
Wir präsentieren den eleganten \textbf{Zykov-Symmetrisierungs}-Beweis.

Nenne einen Graphen \emph{extremal}, wenn er $K_{r+1}$-frei ist und maximale
Kantenzahl hat. Wir behaupten, jeder extremale Graph ist $T(n,r)$.

\textbf{Schritt 1: Extremale Graphen sind vollständig multipartit.}
Sei $G$ extremal. Haben zwei nicht-adjazente Knoten $u, v$ den Grad
$\grad(u) < \grad(v)$, ersetze $u$ durch eine Kopie von $v$ (gleiche Nachbarschaft
wie $v$). Dies erzeugt kein $K_{r+1}$ (da $v$s Nachbarschaft $K_{r+1}$-frei
ist) und kann die Kantenzahl nicht verringern. Also können wir annehmen: für
alle nicht-adjazenten $u,v$ gilt $\grad(u) = \grad(v)$.

Nicht-Adjazenz ist eine Äquivalenzrelation (sind $u \not\sim v$ und
$u \not\sim w$, müssen sie gleichen Grad haben, also gilt $v \not\sim w$
ebenfalls nach Maximalität). Damit ist $G$ vollständig multipartit.

\textbf{Schritt 2: Teile müssen ausgewogen sein.}
Ist $G = K_{n_1, \ldots, n_r}$ mit $n_i > n_j + 1$ für ein $i \ne j$,
verschiebe einen Knoten von Teil $i$ nach Teil $j$. Dies erhöht die Kantenzahl
(die Veränderung ist $n_j - n_i + 1 > 0$ nach sorgfältiger Berechnung),
was der Extremalität widerspricht. Also haben alle Teile Größe
$\lfloor n/r\rfloor$ oder $\lceil n/r\rceil$, d.\,h.\ $G = T(n,r)$.
\end{proof}

\begin{theorem}[Zarankiewicz-Problem, Kővári-Sós-Turán]
\label{satz:zarankiewicz}
Ist $G$ ein bipartiter Graph auf Teilen der Größe $m$ und $n$ (mit $m \le n$),
der kein $K_{s,t}$ als Teilgraph enthält (mit $s \le t$), so gilt:
\[
  |E(G)| \le \tfrac{1}{2}\bigl((t-1)^{1/s} n m^{1-1/s} + (s-1)m\bigr).
\]
Insbesondere gilt ex$(n, K_{2,2}) = O(n^{3/2})$.
\end{theorem}

\begin{proof}[Beweisskizze]
Zähle Paare $(u, \{v_1, v_2, \ldots, v_s\})$, wobei $u$ auf der $m$-Seite
liegt und die $v_i$ auf der $n$-Seite mit allen $v_i$ adjazent zu $u$. Einerseits
ist diese Zahl $\sum_{u} \binom{\grad(u)}{s} \ge m \cdot \binom{|E|/m}{s}$
nach Konvexität. Andererseits gibt es für jede $s$-Menge $\{v_1, \ldots, v_s\}$
auf der $n$-Seite höchstens $t-1$ gemeinsame Nachbarn auf der $m$-Seite
(sonst existiert $K_{s,t}$). Also ist die Zahl höchstens $(t-1)\binom{n}{s}$.
Kombination und Auflösung nach $|E|$ ergibt das Ergebnis.
\end{proof}

\begin{theorem}[Szemerédi-Regularitätslemma]
\label{satz:szemeredi}
Für jedes $\varepsilon > 0$ und jede ganze Zahl $k_0 \ge 1$ existiert ein
$M = M(\varepsilon, k_0)$, sodass die Knotenmenge jedes Graphen $G$ mit
$n \ge M$ Knoten in $k$ Klassen $V_1, \ldots, V_k$ mit $k_0 \le k \le M$
und $|V_1| \le \cdots \le |V_k| \le |V_1| + 1$ partitioniert werden kann,
sodass alle bis auf höchstens $\varepsilon k^2$ der Paare $(V_i, V_j)$
$\varepsilon$-regulär sind.
\end{theorem}

\begin{proof}[Beweisskizze]
Das \textbf{Energie-Inkrement-Argument}: Definiere die \emph{mittlere quadratische
Dichte} (Energie) $q(P) = \sum_{i<j} \frac{|V_i||V_j|}{n^2} d(V_i, V_j)^2$
für eine Partition $P$, wobei $d(V_i, V_j) = e(V_i, V_j)/(|V_i||V_j|)$ die
Kantendichte ist. Es gilt $0 \le q(P) \le 1$.

Hat die aktuelle Partition zu viele irregulä re Paare, verfeinere sie, indem
jede Klasse mithilfe des ``Zeugen'' der Irregularität aufgeteilt wird. Jeder
Verfeinerungsschritt erhöht die Energie um mindestens $\varepsilon^5$. Da
die Energie durch 1 begrenzt ist, terminiert der Prozess nach höchstens
$\varepsilon^{-5}$ Verfeinerungen mit einer regulären Partition.
\end{proof}

\begin{remark}
Das Szemerédi-Regularitätslemma ist ein Eckpfeiler der modernen Kombinatorik.
Es ermöglicht das \emph{Zähl-Lemma} und das \emph{Entfernungs-Lemma} und
war zentral für den Beweis von Szmerédis Satz über arithmetische Progressionen.
Die Schranke $M(\varepsilon)$ wächst wie ein Exponentialturm der Höhe
$O(1/\varepsilon^5)$, was bekanntermaßen im Wesentlichen optimal ist
(Gowers, 1997).
\end{remark}

%% =============================================================
\section{Literatur}
\label{sec:literatur}
%% =============================================================

\begin{thebibliography}{99}

\bibitem{euler1736}
L.~Euler,
\emph{Solutio problematis ad geometriam situs pertinentis},
Commentarii Academiae Scientiarum Imperialis Petropolitanae \textbf{8} (1736),
128--140.

\bibitem{cayley1889}
A.~Cayley,
\emph{A theorem on trees},
Quart.\ J.\ Math.\ \textbf{23} (1889), 376--378.

\bibitem{appelhaken1977}
K.~Appel und W.~Haken,
\emph{Every planar map is four colorable. Part~I: Discharging},
Illinois J.\ Math.\ \textbf{21} (1977), 429--490.
K.~Appel, W.~Haken und J.~Koch,
\emph{Every planar map is four colorable. Part~II: Reducibility},
Illinois J.\ Math.\ \textbf{21} (1977), 491--567.

\bibitem{robertson1997}
N.~Robertson, D.~P.~Sanders, P.~D.~Seymour und R.~Thomas,
\emph{The four-colour theorem},
J.\ Combin.\ Theory Ser.\ B \textbf{70} (1997), 2--44.

\bibitem{degrey2018}
A.~D.~N.~J.\ de~Grey,
\emph{The chromatic number of the plane is at least 5},
Geombinatorics \textbf{28} (2018/2019), 18--31.
arXiv:1804.02385.

\bibitem{diestel2017}
R.~Diestel,
\emph{Graphentheorie},
5.~Aufl., Springer, Berlin, 2017.

\bibitem{bollobas1998}
B.~Bollob\'{a}s,
\emph{Modern Graph Theory},
Springer, New York, 1998.

\bibitem{karp1972}
R.~M.~Karp,
\emph{Reducibility among combinatorial problems},
in: R.~E.~Miller und J.~W.~Thatcher (Hrsg.),
\emph{Complexity of Computer Computations},
Plenum Press, New York, 1972, S.~85--103.

\bibitem{turan1941}
P.~Tur\'{a}n,
\emph{On an extremal problem in graph theory} (auf Ungarisch),
Mat.\ Fiz.\ Lapok \textbf{48} (1941), 436--452.

\bibitem{szemeredi1978}
E.~Szemer\'{e}di,
\emph{Regular partitions of graphs},
Probl\`{e}mes Combinatoires et Th\'{e}orie des Graphes,
Colloques Internationaux du CNRS \textbf{260}, Paris, 1978, S.~399--401.

\bibitem{mosermoser1961}
L.~Moser und W.~Moser,
\emph{Solution to problem 10},
Canad.\ Math.\ Bull.\ \textbf{4} (1961), 187--189.

\bibitem{kuratowski1930}
K.~Kuratowski,
\emph{Sur le probl\`{e}me des courbes gauches en topologie},
Fund.\ Math.\ \textbf{15} (1930), 271--283.

\end{thebibliography}

\end{document}
